% Header aus der Vorlage
\documentclass[a4paper,11pt, footheight=26pt]{scrartcl}
\usepackage[head=23pt]{geometry}	% head=23pt umgeht Fehlerwarnung, dafür größeres "top" in geometry
\geometry{a4paper, top=30mm, left=20mm, right=20mm, bottom=22mm,headsep=10mm, footskip=12mm}

% Input inkl. Umlaute, Silbentrennung
\usepackage[T1]{fontenc}
\usepackage[utf8]{inputenc}
\usepackage[ngerman]{babel}
\usepackage{csquotes}	% Anführungszeichen

% HTW Corporate Design: Arial (Helvetica)
\usepackage{helvet}
\renewcommand{\familydefault}{\sfdefault}

% Style-Aufhübschung
\usepackage{soul, color}	% Kapitälchen, Unterstrichen, Durchgestrichen usw. im Text
\usepackage{scrlayer-scrpage}	% Kopf-/Fußzeile
%\usepackage{titleref}
\usepackage[perpage]{footmisc}	% Fußnotenzählung Seitenweit, nicht Dokumentenweit

% Mathe usw.
\usepackage{amssymb}
\usepackage[fleqn]{amsmath}	% fleqn: align-Umgebung rechtsbündig
\usepackage{xcolor}
\usepackage{esint}	% Schönere Integrale, \oiint vorhanden
\everymath=\expandafter{\the\everymath\displaystyle}	% Mathe Inhalte werden weniger verkleinert
\usepackage{wasysym}	% mehr Symbole, bspw \lightning

% tikz usw.
\usepackage{tikz}
\usepackage{pgfplots}
\pgfplotsset{compat=1.11}	% Umgeht Fehlermeldung
\usetikzlibrary{graphs}
%\usetikzlibrary{through}	% ???
\usetikzlibrary{arrows}
\usetikzlibrary{arrows.meta}	% Pfeile verändern / vergrößern: \draw[-{>[scale=1.5]}] (-3,5) -> (-3,3);
\usetikzlibrary{automata,positioning} % Zeilenumbruch im Node node[align=center] {Text\\nächste Zeile}
\usetikzlibrary{patterns}	% Schraffierte Füllung
\tikzstyle{reverseclip}=[insert path={	% Inverser Clip \clip
	(current page.north east) --
	(current page.south east) --
	(current page.south west) --
	(current page.north west) --
	(current page.north east)}
% Nutzen: 
%\begin{tikzpicture}[remember picture]
%\begin{scope}
%\begin{pgfinterruptboundingbox}
%\draw [clip] DIE FLÄCHE, IN DER OBJEKT NICHT ERSCHEINEN SOLL [reverseclip];
%\end{pgfinterruptboundingbox}
%\draw DAS OBJEKT;
%\end{scope}
%\end{tikzpicture}
]	% Achtung: dafür muss doppelt kompliert werden!
\usepackage{graphpap}	% Grid für Graphen

% Tabular
\usepackage{longtable}	% Große Tabellen über mehrere Seiten
\usepackage{multirow}	% Multirow/-column: \multirow{2[Anzahl der Zeilen]}{*[Format]}{Test[Inhalt]} oder \multicolumn{7[Anzahl der Reihen]}{|c|[Format]}{Test2[Inhalt]}
\renewcommand{\arraystretch}{1.3} % Tabellenlinien nicht zu dicht
\usepackage{colortbl}
\arrayrulecolor{gray}	% heller Tabellenlinien

% Nützliches
\usepackage{verbatim}	% u.a. zum auskommentieren via \begin{comment} \end{comment}
\usepackage{tabto}	% Tabs: /tab zum nächsten Tab oder /tabto{.5 \CurrentLineWidth} zur Stelle in der Linie
\NumTabs{6}	% Anzahl von Tabs pro Zeile zum springen
\usepackage{listings} % Source-Code mit Tabs
\usepackage{lstautogobble} 
\usepackage{enumitem}	% Anpassung der enumerates
\setlist[enumerate,1]{label=\arabic*.)}	% global andere Enum-Items
\usepackage{letltxmacro} % neue Definiton von Grundbefehlen
% Nutzen:
%\LetLtxMacro{\oldemph}{\emph}
%\renewcommand{\emph}[1]{\oldemph{#1}}

% Einrichtung von lst
\lstset{
basicstyle=\ttfamily, 
mathescape=true, 
escapeinside=||, 
autogobble, 
tabsize=2,
basicstyle=\footnotesize\sffamily\color{black},
frame=single,
rulecolor=\color{lightgray},
numbers=left,
numbersep=5pt,
%basicstyle=\footnotesize\sffamily\color{black},
commentstyle=\color{gray},
%frame=single,
%numbers=left,
%numbersep=5pt,
numberstyle=\tiny\color{gray},
keywordstyle=\color{green},
%showspaces=false,
showstringspaces=false,
stringstyle=\color{orange},
tabsize=2,
literate=%
    {Ö}{{\"O}}1
    {Ä}{{\"A}}1
    {Ü}{{\"U}}1
    {ß}{{\ss}}1
    {ü}{{\"u}}1
    {ä}{{\"a}}1
    {ö}{{\"o}}1
    {~}{{\textasciitilde}}1
}

% BibTeX
\usepackage[backend=bibtex, bibencoding=ascii]{biblatex}	% BibTeX
\usepackage{makeidx}
%\makeglossary
%\makeindex

% Grafiken
\usepackage{graphicx}
\usepackage{epstopdf}	% eps-Vektorgrafiken einfügen

% pdf-Setup
\usepackage{pdfpages}
\usepackage[bookmarks,%
bookmarksopen=false,% Klappt die Bookmarks in Acrobat aus
colorlinks=true,%
linkcolor=black,%
citecolor=red,%
urlcolor=green,%
]{hyperref}

% Titel, Autor usw. werden vor dem Anfang des Dokuments in einem Rutsch definiert…
\newcommand{\DTitel}[1]{\newcommand{\Dokumententitel}{#1}}
\newcommand{\DUntertitel}[1]{\newcommand{\Dokumentenuntertitel}{#1}}
\newcommand{\DAutor}[1]{\newcommand{\Dokumentenautor}{#1}}
\newcommand{\DNotiz}[1]{\newcommand{\Dokumentennotiz}{#1}}
% … Deswegen folgendes erst Nach Dokumentenbeginn ausführen:
\AtBeginDocument{
	\hypersetup{
		pdfauthor={\Dokumentenautor},
		pdftitle={HTW Dresden | \Dokumententitel - \Dokumentenuntertitel},
	}
	\automark[section]{section}
	\automark*[subsection]{subsection}
	\pagestyle{scrheadings}
	\ihead{\includegraphics[height=1.7em]{../LaTeX_master/HTW-Logo.eps}}
	\ohead{\Dokumententitel}
	\ofoot{\footnotesize{\textcolor{darkgray}{Mitschrift von\\ \Dokumentenautor}}}
	% Titelseite
	\title{\includegraphics[width=0.35\textwidth]{../LaTeX_master/HTW-Logo.eps}\\\vspace{0.5em}
	\Huge\textbf{\Dokumententitel} \\\vspace*{0,5cm}
	\Large \Dokumentenuntertitel \\\vspace*{4cm}}
	\author{\textcolor{darkgray}{Mitschrift von \Dokumentenautor} \vspace*{1cm}\\\Dokumentennotiz}
}

%% EIGENE BEFEHLE

%Farbdefinitionen
\definecolor{red}{RGB}{180,0,0}
\definecolor{green}{RGB}{75,160,0}
\definecolor{blue}{RGB}{0,75,200}
\definecolor{orange}{RGB}{255,128,0}
\definecolor{yellow}{RGB}{255,245,0}
\definecolor{purple}{RGB}{75,0,160}
\definecolor{cyan}{RGB}{0,160,160}
\definecolor{brown}{RGB}{120,60,10}

% Textfarbe ändern
\newcommand{\tred}[1]{\textcolor{red}{#1}}
\newcommand{\tgreen}[1]{\textcolor{green}{#1}}
\newcommand{\tblue}[1]{\textcolor{blue}{#1}}
\newcommand{\torange}[1]{\textcolor{orange}{#1}}
\newcommand{\tyellow}[1]{\textcolor{yellow}{#1}}
\newcommand{\tpurple}[1]{\textcolor{purple}{#1}}
\newcommand{\tcyan}[1]{\textcolor{cyan}{#1}}
\newcommand{\tbrown}[1]{\textcolor{brown}{#1}}

% Umstellen der Tabellen Definition
\newcommand{\mpb}[1][.3]{\begin{minipage}{#1\textwidth}\vspace*{3pt}}
\newcommand{\mpe}{\vspace*{3pt}\end{minipage}}

\newcommand{\resultul}[1]{\underline{\underline{#1}}}
\newcommand{\parskp}{$ $\\}	% new line after paragraph
\newcommand{\corr}{\;\widehat{=}\;}
\newcommand{\mdeg}{^{\circ}}

\newcommand{\nok}[2]{\begin{pmatrix}#1\\#2\end{pmatrix}}	% n über k

\bibliography{../Literatur/HTW_Literatur.bib}

% Definition von Titel, Autor usw.
\DTitel{Elektronik}
\DUntertitel{Übungsskript}
\DAutor{Falk Jonatan Strube}
\DNotiz{Übung von Prof. Dr.-Ing. Flach (bis 12/2015)}

\begin{document}

\pagenumbering{Roman}

\maketitle
\newpage
\tableofcontents
\newpage

\pagenumbering{arabic}

\section{Übung 1}

\section*{1.2}

$\boxed{P=M\cdot 2\pi \cdot n}$\\
P … mech. Leistung\\
M … Drehmoment\\
n … Drezahl
\paragraph{a)} ges.: $P$ für\\
$M_1= 1Nm$, $n=1000 min^{-1}$\\
$M_2= 2Nm$, $n=1000 min^{-1}$\vspace*{1em}\\
$P_1=1\underset{Ws}{Nm} \cdot 2\pi \cdot 1000 \underset{\frac{1}{60}s^{-1}}{min^{-1}}=\resultul{104,72W}$\\
$P_2=\resultul{209,44W}$

\paragraph{b)} zugeschnittene Größengleichung für $P$, $M$ in $Nm$, $n$ in $min^{-1}$, $P$ in $Watt$\\
$\frac{P}{W}=\frac{M}{Nm} \cdot 2 \pi \cdot \frac{n}{\underbrace{s^{-1}}_{60\cdot min^{-1}}}=2\pi \cdot \frac{1}{60} \cdot \frac{M}{Nm} \cdot \frac{n}{min^{-1}}=\underset{\text{zugeschnittene Größengleichung}}{\boxed{0,105\frac{M}{Nm}\cdot \frac{n}{min^{-1}}=\frac{P}{W}}}$

\section*{1.3}
$\boxed{W=P\cdot t}$
\paragraph{a)} Zugeschnitten Größengleichung für $W$ in $kWh$, $P$ in $kW$, $t$ in $s$\\
$\frac{W}{kWh}=\frac{P}{kW}\cdot \frac{t}{\underbrace{h}_{3600s}}=\frac{1}{3600}\cdot \frac{P}{kW} \cdot \frac{t}{s}$

\paragraph{b)} ges.: $W$ für\\
$P_1=0,4kW$, $t=7200s$\\
$P_2=1,2kW$, $t=3600s$\medskip\\
$W_1=\frac{1}{3600} \cdot 0,4 \cdot 7200kWh=\resultul{0,8kWh}$\smallskip\\
$W_2=\resultul{1,2kWh}$

\section*{1.4}
Luftgeschwindigkeit\\
Staurohr nach Prandl: $\Delta p = (\frac{\varrho}{2})v^2$ mit $\varrho$ … Luftdichte $1,1 \frac{kg}{m^3}$\\
ges.: $v$ in $\frac{km}{h}$ bei $\varrho$ in $\frac{kg}{m^3}$, $\Delta p$ in $Pa \Rightarrow 1 Pa = 1\frac{N}{m^2}= 1 \frac{kg}{s^2}$\\
$v=\sqrt{\frac{2\cdot \Delta p}{\varrho}}=\sqrt{\frac{2\Delta \frac{p}{\frac{kg}{m\cdot s^2}}}{\frac{\varrho}{\frac{kg}{m^3}}}}=\sqrt{\frac{\frac{2\Delta p}{\varrho}}{\frac{m^2}{s^2}}}=\sqrt{\frac{\frac{2 \Delta p}{\varrho}}{(\frac{10^{-3}\cdot km}{\frac{1}{3600}h})^2}}=\sqrt{\frac{2}{3,6^2\cdot 1,1}}\sqrt{\Delta p}\frac{km}{h}$\\
$v=\resultul{0,37\sqrt{\Delta p}\frac{km}{h}}$\medskip\\
ges.: $v_i$\\
$\Delta p = 10,30,50,100 Pa$\\
$v_1=\resultul{0,37\cdot \sqrt{10}\frac{km}{h}=1,17\frac{km}{h}}$\\
$v_2=\resultul{2,03\frac{km}{h}}$\\
$v_3=\resultul{2,62\frac{km}{h}}$\\
$v_4=\resultul{3,7 \frac{km}{h}}$

\section*{2.1}

geg.: $I=2mA$, $t=2 min$\\
ges.: $Q$\\
Lösung: $\boxed{I=\frac{\Delta Q}{\Delta t}=\frac{Q}{t}}$\\
$Q=I\cdot t=2mA \cdot 2 \cdot 60s = 240 mAs = \resultul{0,24 As}$

\section*{2.2}

geg.: $I=125mA$, $t=3,2s$\\
ges.: Anzahl Elektronen $n$\\
$e = 1,602 \cdot 10^{-19}C$\\
Lösung: $Q=I\cdot t=n\cdot e$\\
$n=\frac{I\cdot t}{e}=\frac{125mA \cdot 3,2s}{1,602\cdot 10^{-19}As}=\resultul{250\cdot 10^{16}\; Elektronen}$

\section*{2.3}

geg.: $n_p=0,2\cdot 10^{11}$, $n_n=10^9$, $e=1,602\cdot 10^{-10} C$, $t=1s$\\
ges.: $I$\\
Lösung: $I=\frac{\Delta Q}{\Delta t}=\frac{e(n_n+n_p)}{\Delta p}=\frac{1,602\cdot 10^{-19}As(26\cdot 10^9}{1s}=\resultul{41,652 \cdot 10^{-10}A}$

\section*{2.4}

geg.: $Q=20mAh$, $I=1 \mu A$\\
ges.: $t$, $n$\\
Lösung: $t=\frac{Q}{I}=\frac{20\cdot 10^{-3}Ah}{1 \cdot 10^{-6}A}=\resultul{20000h}$\\
$n=\frac{Q}{e}=\frac{20\cdot 10^{-3}\cdot 3600 As}{1,602\cdot 10^{-19} As}=\resultul{4,944\cdot 10^{20}}$

\section*{2.5}
Die Stromstärke beim Laden und Entladen einer Autobatterie hat folgenden Zeitverlauf:\\
\begin{tabular}{l | l}
Zeitintervall & Stromstärke \\
\hline
$0 \leq t < 3h$ & $2A$\\
$3 \leq t < 4h$ & $-0,5A$\\
$t\geq 4h$ & $0A$
\end{tabular}\\
Welche Ladung hat die Batterie nach $4h$, wenn Anfangsladung $0$ war?\\
Zeitverlauf von $I$ und $Q$ ist darzustellen.

\paragraph{Lösung:} $Q=2A \cdot 3h + (-0,5 A) \cdot 1 h = 6Ah - 0,5Ah = \resultul{5,5Ah}$\\
\begin{tikzpicture}
% Achsen zeichnen, zuerst hier, da clip die Zeichenflaeche beschneidet
\draw[->,thick] (0,0) -- (4.5,0) node[right] {$\frac{t}{h}$};
\draw[->,thick] (0,-.75) -- (0,2.8) node[above right] {$\frac{I}{A},\textcolor{cyan}{\frac{Q}{Ah}}$};
% Achsen beschriften
\foreach \x in {1,2,3,4}
\draw (\x,-.1) -- (\x,.1) node[below=5pt] {$\scriptstyle\x$};
\foreach \y in {-0.5,1,2}
\draw (-.1,\y) -- (.1,\y) node[left=5pt] {$\scriptstyle\y$};
\draw[cyan] (-.1,2) -- (.1,2) node[left=10pt] {$\scriptstyle6$};

\draw[gray, thick] (0,2) -- (3,2) -- (3,-0.5) -- (4,-0.5) -- (4,0);
\draw[cyan, thick] (0,0) -- (3,2) -- (4,1.83);
\end{tikzpicture}






\newpage
%\printbibliography
\end{document}